\section{Native Functional Coverage}\label{sec:coverage}
All seven scoped teachers have passing packages for each of the three
architectures: \textbf{21/21 combinations}, five selected checkpoints per
combination and 1,000/1,000 proper complete invariant-valid executions per
checkpoint. The 105 policies account for 105,000 successful observations.
This establishes broad finite functionality within supplied interpreters;
it does not establish universal correctness or equivalence of architectures.

\begin{table}[htbp]\centering\small
\begin{tabular}{lccc}
\toprule Teacher & E1 & DeepSets & Transformer \\
\midrule
Luby & H: 5/5 & H: 5/5 & H: 5/5 \\
Johansson & H: 5/5 & H: 5/5 & H: 5/5 \\
ColorTrials & N: 5/5 & N: 5/5 & N: 5/5 \\
JP & H: 5/5 & H: 5/5 & H: 5/5 \\
BE base & N: 5/5 & N: 5/5 & N: 5/5 \\
GM & N: 5/5 & N: 5/5 & N: 5/5 \\
Greedy & H: 5/5 & H: 5/5 & H: 5/5 \\
\bottomrule
\end{tabular}
\caption{Native functional coverage. Each passing checkpoint completes all
1,000 runs. H: twelve historical passive packages; N: nine later CF packages.
The cohorts have different training and final-test protocols and do not
constitute a homogeneous architecture-ranking experiment.}
\label{tab:coverage}
\end{table}

Historical packages use seeds 83101--83105; later ColorTrials/BE/GM packages
use 94101--94105. Representative BE/ColorTrials policies use 4,000 updates,
GM 500, while acquired-label budgets vary by package. Matched controls also
pass all nine later combinations: CF is not necessary for native functional
success in these tests. Final sizes are 32,48,64; seven historical versus
ten later graph strata have different weights. Appendix~\ref{app:coverage}
retains confidence-limit assumptions, unsuccessful candidates and the full
budget/provenance table. The following continuation analyses concern
separate historical or prospectively frozen cohorts; the 21/21 result is
not attributed to every model in those analyses.
